\cleardoublepage
\section*{Tóm tắt}
\addcontentsline{toc}{section}{Tóm tắt}
Luận văn này tập trung nghiên cứu và xây dựng một ứng dụng web hỗ trợ người dùng trong việc
tìm kiếm thông tin quán ăn, tạo và quản lý các hành trình ăn uống (food tour), đồng thời xác định
lộ trình di chuyển hợp lý giữa nhiều địa điểm ẩm thực trong cùng một hành trình.

Hệ thống cho phép người dùng tìm kiếm và xem thông tin chi tiết về các quán ăn, thêm một hoặc
nhiều quán ăn vào food tour cá nhân và sắp xếp thứ tự ghé thăm nhằm tối ưu hóa quãng đường
di chuyển. Bên cạnh đó, luận văn đề xuất cơ chế gợi ý food tour dựa trên mức độ tương đồng
giữa các quán ăn trong lịch sử sử dụng của người dùng và các food tour đã tồn tại, giúp người
dùng dễ dàng tham khảo và lựa chọn các hành trình phù hợp.

Ngoài ra, hệ thống còn hỗ trợ chức năng chia sẻ food tour ở chế độ công khai, cho phép người
dùng khác tìm kiếm và sử dụng lại các food tour đã được chia sẻ. Luận văn trình bày chi tiết quá
trình phân tích yêu cầu, thiết kế hệ thống, xây dựng cơ sở dữ liệu, kiến trúc tổng thể và triển khai
ứng dụng. Kết quả đạt được cho thấy hệ thống đáp ứng tốt các yêu cầu chức năng đã đề ra và có
tiềm năng mở rộng trong tương lai.
\cleardoublepage
\section*{Abstract}
\addcontentsline{toc}{section}{Abstract}

This thesis focuses on the design and development of a web-based application that supports
users in searching for restaurant information, creating and managing food tours, and determining
efficient travel routes among multiple dining locations within a single itinerary.

The proposed system allows users to search for and view detailed information about restaurants,
add one or multiple restaurants to a personal food tour, and arrange the visiting order in order to
optimize travel distance and time. In addition, the thesis proposes a food tour recommendation
mechanism based on the similarity of restaurants appearing in users’ historical food tours and
existing food tours, enabling users to easily discover and select suitable itineraries.

Furthermore, the system provides a public sharing feature that allows users to publish their food
tours so that other users can search for and reuse shared food tours. The thesis presents the
complete process of requirement analysis, system design, database design, overall architecture,
and application implementation. The experimental results demonstrate that the system satisfies
the defined functional requirements and shows potential for further extension in the future.